% Generated draft: requires mathematical and source review.
\documentclass[11pt]{article}
\usepackage[margin=25mm]{geometry}
\usepackage{fontspec}
\setmainfont{Latin Modern Roman}
\usepackage{amsmath,amssymb,mathtools,unicode-math}
\setmathfont{Latin Modern Math}
\usepackage{enumitem,needspace,xurl,hyperref}
\setlength{\parindent}{0pt}
\setlength{\parskip}{5pt}
\setlength{\emergencystretch}{3em}
\newcommand{\N}{\mathbb{N}}
\newcommand{\Z}{\mathbb{Z}}
\newcommand{\Q}{\mathbb{Q}}
\newcommand{\R}{\mathbb{R}}
\newcommand{\C}{\mathbb{C}}
\newcommand{\F}{\mathbb{F}}
\newcommand{\A}{\mathbb{A}}
\newcommand{\PP}{\mathbb{P}}
\newcommand{\E}{\mathbb{E}}
\newcommand{\eps}{\varepsilon}
\newcommand{\dd}{\,\mathrm{d}}
\begin{document}
\section*{UnsolvedMath ID: 30006495}
\textbf{Source problem code:} catalogue entry
\subsection{Definitions and mathematical statement}
Let \(\mathcal C\) denote a braided tensor category, with monoidal unit \(\mathbf 1\) and braiding \(c_{X,Y}:X\otimes Y\to Y\otimes X\). The source does not specify the base field, linearity, finiteness, exactness, or existence assumptions on \(\mathcal C\).

To explain the usual meaning of a Nichols algebra object, suppose the following constructions exist in \(\mathcal C\). For an object \(V\), form the graded tensor algebra
\[
T_{\mathcal C}(V)=\bigoplus_{n\geq 0}V^{\otimes n},\qquad V^{\otimes 0}=\mathbf 1,
\]
with its braided Hopf algebra structure determined by declaring \(V\) primitive. If there is a largest homogeneous braided Hopf ideal \(I\) with zero components in degrees \(0\) and \(1\), define
\[
\mathcal B_{\mathcal C}(V)=T_{\mathcal C}(V)/I.
\]
Existence of the direct sum, this Hopf algebra structure, the largest ideal, and the quotient is a qualification of this definition; it is not established for every category covered by the source's word ``arbitrary.'' The source also does not specify whether its ``generalized Nichols algebras'' are precisely these objects or a broader class.

\textbf{Source research problem.} Develop a root-system theory and a Weyl-groupoid theory for Nichols algebra objects in the intended more general braided tensor categories, and classify the resulting generalized Nichols algebras. The supplied normalized statement requests this for arbitrary braided tensor categories. The original excerpt instead asks for root-system theory and classification ``in such more general settings,'' without specifying a universal domain or explicitly mentioning Weyl groupoids.

The motivating exposition concerns braided tensor categories with large commutative algebra objects that, under unspecified conditions, resemble Drinfeld centers associated with Nichols algebras. It suggests categories related to affine Lie algebras at integer level and to other \(W\)-algebras, including settings involving inverse Hamiltonian reduction. Neither the relevant categories nor the meaning of ``large'' nor the precise center construction is supplied.

\textbf{Unresolved formulation.} No axioms for the requested root systems, objects and reflection morphisms of the Weyl groupoid, admissible class of algebra objects, finiteness restrictions, or equivalence relation for classification are specified. Thus the record determines a research program, but not a unique fully quantified conjecture. The conditional definition above explains standard terminology without supplying the missing scope or hypotheses.
\subsection{Short English statement}
The problem asks for root systems, Weyl groupoids, and a classification for Nichols algebras in more general braided tensor categories. The supplied record leaves the categories, existence hypotheses, classification criteria, and meaning of ``generalized'' unspecified. Its dated partially solved assessment concerns the logarithmic Kazhdan--Lusztig correspondence, whose relationship to this classification problem needs review.
\subsection{Sources}
\begin{itemize}
\item Supplied local registry record: permanent UnsolvedMath ID \(30006495\), title ``Root Systems and Classification of Generalized Nichols Algebras,'' exact source problem code \texttt{catalogue entry}. The proposer is not identified; the supplied proposed year is 2025.
\item Supplied report citation: \emph{Mini-Workshop: The Yang--Baxter Equation and Representations of Braid Groups} (2026), \emph{Oberwolfach Reports}, pp.~2633--2672, DOI \texttt{10.4171/owr/2025/49}. \url{https://doi.org/10.4171/owr/2025/49}. The supplied excerpt ends with a report header and page number 2657. The citation's 2026 date differs from the 2025 proposed year, report label, and DOI year; this discrepancy is unresolved.
\item Supplied literature source: ``Conditional algebraic proof of the logarithmic Kazhdan--Lusztig correspondence.'' \url{https://arxiv.org/abs/2501.10735}. No author attribution is provided in the record.
\item Supplied literature source: ``Oberwolfach Report 49/2025.'' \url{https://doi.org/10.4171/owr/2025/49}.
\item Literature checked at: \texttt{2026-08-22}. This is supplied dated curation, not an independently verified assessment. The supplied status is ``partially solved.'' The assessment states that the logarithmic Kazhdan--Lusztig correspondence has a general algebraic proof under explicit analytic and categorical hypotheses, while the unconditional analytic hypotheses are not established in full generality. Those hypotheses are not included in the record, and this assessment does not by itself establish partial resolution of the root-system and classification problem.
\item The registry labels the statement \texttt{corrected\_verified} and reports that the official report was checked, the question isolated, and notation repaired. These are supplied verification claims; no independent source consultation or mathematical verification is asserted here.
\end{itemize}
\end{document}
